% p.76 — الباب الثالث والثلاثين في تمرين الصرف والتركيب والاعراب في الآيات
% Layout: \bab heading across the page, then twocol. rightcol: الاستعاذة, basmala, الفاتحة,
% basmala, البقرة 1–5, basmala, الناس; leftcol: basmala, الفلق, basmala, الاخلاص, basmala,
% الكافرون, basmala, الكوثر, closing du'ā'. Items are \par-separated. The lithograph marks the
% end of a basmala / sura with a small circle; it is rendered here as {\latinfont°}.
% Vocalisation (including sukūn on letters of prolongation and dagger alifs) as printed.
\mspage{76}
\bab{الباب الثالث والثلاثين في تمرين الصرف والتركيب والاعراب في الآيات}
\begin{twocol}
\begin{rightcol}
اَعُوْذُ بِاللهِ مِنَ الشَّيْطَانِ الرَّجِيْمِ
\par
بِسْمِ اللهِ الرَّحْمٰنِ الرَّحِيْمِ {\latinfont°}
\par
اَلْحَمْدُ لِلّٰهِ رَبِّ الْعَالَمِيْنَ الرَّحْمٰنِ الرَّحِيْمِ مَالِكِ يَوْمِ الدِّيْنِ اِيَّاكَ نَعْبُدُ وَاِيَّاكَ نَسْتَعِيْنُ اِهْدِنَا الصِّرَاطَ الْمُسْتَقِيْمَ صِرَاطَ الَّذِيْنَ اَنْعَمْتَ عَلَيْهِمْ غَيْرِ الْمَغْضُوْبِ عَلَيْهِمْ وَلَا الضَّالِّيْنَ {\latinfont°}
\par
بِسْمِ اللهِ الرَّحْمٰنِ الرَّحِيْمِ {\latinfont°}
\par
الٓمّٓ ذَالِكَ الْكِتٰبُ لَا رَيْبَ فِيْهِ هُدًى لِلْمُتَّقِيْنَ الَّذِيْنَ يُؤْمِنُوْنَ بِالْغَيْبِ وَيُقِيْمُوْنَ الصَّلٰوةَ وَمِمَّا رَزَقْنَاهُمْ يُنْفِقُوْنَ وَالَّذِيْنَ يُؤْمِنُوْنَ بِمَا اُنْزِلَ اِلَيْكَ وَمَا اُنْزِلَ مِنْ قَبْلِكَ وَبِالْاٰخِرَةِ هُمْ يُوْقِنُوْنَ اُولٰئِكَ عَلٰى هُدًى مِنْ رَبِّهِمْ وَاُولٰئِكَ هُمُ الْمُفْلِحُوْنَ
\par
بِسْمِ اللهِ الرَّحْمٰنِ الرَّحِيْمِ {\latinfont°}
\par
قُلْ اَعُوْذُ بِرَبِّ النَّاسِ مَلِكِ النَّاسِ اِلٰهِ النَّاسِ مِنْ شَرِّ الْوَسْوَاسِ الْخَنَّاسِ الَّذِيْ يُوَسْوِسُ فِيْ صُدُوْرِ النَّاسِ مِنَ الْجِنَّةِ وَالنَّاسِ {\latinfont°}
\par
\end{rightcol}
\begin{leftcol}
بِسْمِ اللهِ الرَّحْمٰنِ الرَّحِيْمِ
\par
قُلْ اَعُوْذُ بِرَبِّ الْفَلَقِ مِنْ شَرِّ مَا خَلَقَ وَمِنْ شَرِّ غَاسِقٍ اِذَا وَقَبَ وَمِنْ شَرِّ النَّفّٰثٰتِ فِي الْعُقَدِ وَمِنْ شَرِّ حَاسِدٍ اِذَا حَسَدَ
\par
بِسْمِ اللهِ الرَّحْمٰنِ الرَّحِيْمِ {\latinfont°}
\par
قُلْ هُوَ اللهُ اَحَدٌ اَللهُ الصَّمَدُ لَمْ يَلِدْ وَلَمْ يُوْلَدْ وَلَمْ يَكُنْ لَهُ كُفُوًا اَحَدٌ .
\par
بِسْمِ اللهِ الرَّحْمٰنِ الرَّحِيْمِ {\latinfont°}
\par
قُلْ يَا اَيُّهَا الْكَافِرُوْنَ لَا اَعْبُدُ مَا تَعْبُدُوْنَ وَلَا اَنْتُمْ عَابِدُوْنَ مَا اَعْبُدُ وَلَا اَنَا عَابِدٌ مَا عَبَدْتُمْ وَلَا اَنْتُمْ عَابِدُوْنَ مَا اَعْبُدُ لَكُمْ دِيْنُكُمْ وَلِيَ دِيْنِ {\latinfont°}
\par
بِسْمِ اللهِ الرَّحْمٰنِ الرَّحِيْمِ {\latinfont°}
\par
اِنَّا اَعْطَيْنَاكَ الْكَوْثَرَ فَصَلِّ لِرَبِّكَ وَانْحَرْ اِنَّ شَانِئَكَ هُوَ الْاَبْتَرُ {\latinfont°}
\par
رَبَّنَا اٰتِنَا فِي الدُّنْيَا حَسَنَةً وَفِي الْاٰخِرَةِ حَسَنَةً وَقِنَا عَذَابَ النَّارِ رَبَّنَا ظَلَمْنَا اَنْفُسَنَا وَاِنْ لَمْ تَغْفِرْ لَنَا وَتَرْحَمْنَا لَنَكُوْنَنَّ مِنَ الْخَاسِرِيْنَ اِنَّا نَسْئَلُكَ عِلْمًا نَافِعًا وَعَمَلًا مُتَقَبَّلًا وَرِزْقًا حَلَالًا طَيِّبًا وَشِفَاءً مِنْ كُلِّ دَاءٍ . وَالْحَمْدُ لِلّٰهِ رَبِّ الْعَالَمِيْنَ
\par
\end{leftcol}
\end{twocol}
